\documentclass[11pt]{article}
\usepackage[margin=2.4cm]{geometry}
\usepackage{amsmath,amssymb,bm}
\usepackage{xcolor}
\usepackage{hyperref}
\hypersetup{colorlinks=true,linkcolor=blue!50!black,urlcolor=blue!50!black}

\newcommand{\us}{u_s}
\newcommand{\un}{u_n}
\newcommand{\Us}{\bar U_s}
\newcommand{\ds}{\partial_s}
\newcommand{\dn}{\partial_n}
\newcommand{\dt}{\partial_t}
\newcommand{\code}[1]{\texttt{\small #1}}

\title{\bf Derivation of every equation in \code{sw\_sn\_driver.py}\\
\large Full shallow water on a finite-amplitude meander, channel-fitted $(s,n)$}
\author{\code{numerical/dedalus\_meander\_full\_SW}}
\date{}

\begin{document}
\maketitle
\vspace{-8mm}

\noindent This note derives, in order, every formula that appears in
\code{sw\_sn\_driver.py}. Each section names the code symbol it produces.

\tableofcontents

% ===================================================================
\section{Why $(s,n)$, and the metric $\sigma=1+nC$}
\label{sec:metric}

The channel meanders with \emph{finite} amplitude, so its banks are not straight
lines and a Chebyshev grid cannot be laid on them directly. We therefore use
\textbf{channel-fitted} coordinates: $s$ = arc length along the centreline,
$n$ = signed distance across it. In $(s,n)$ the channel is a straight strip
$n\in[-b,b]$ whatever the planform does.

Let the centreline be $\bm r_c(s)$ with unit tangent $\hat{\bm t}(s)$ and unit
normal $\hat{\bm n}(s)$, and let $C(s)$ be its curvature. A point of the domain is
\begin{equation}
  \bm r(s,n)=\bm r_c(s)+n\,\hat{\bm n}(s).
\end{equation}
With the Frenet relation $\mathrm d\hat{\bm n}/\mathrm ds=+C\,\hat{\bm t}$
(this \emph{defines} our sign convention: $n$ increases \emph{away} from the centre
of curvature; equivalently the direction angle is $\theta(s)=-\int^s C\,\mathrm ds'$),
\begin{equation}
  \frac{\partial\bm r}{\partial s}=\hat{\bm t}+n\frac{\mathrm d\hat{\bm n}}{\mathrm ds}
  =(1+nC)\,\hat{\bm t},
  \qquad
  \frac{\partial\bm r}{\partial n}=\hat{\bm n},
\end{equation}
so the coordinates are orthogonal with scale factors
\begin{equation}
  \boxed{\;h_s=\sigma\equiv 1+nC(s),\qquad h_n=1.\;}
  \tag{$\to$ \code{sigma\_metric}, \code{sig}}
\end{equation}
$\sigma$ is the local stretching of an along-channel line at offset $n$: the outer
bank ($n$ with the sign of $C$) is longer, the inner bank shorter.

\paragraph{Validity.} $\sigma>0$ is required or the mapping folds. Since
$|n|\le b$, this is $|C|b<1$: the bend radius must exceed the half-width. The test
\code{tests/test\_base\_profiles.py} asserts $\min\sigma>0$.

\paragraph{Sign convention (important).} Mirroring the channel ($C\to-C$) maps the
problem to itself with $n\to-n$, so the linear growth rate $\sigma(k)$ and phase
speed $c(k)$ are \emph{unchanged}; only the labelling of ``outer'' vs ``inner''
bank flips. The convention above is the one used by the solver
($\sigma=1+nC$); the lab-frame renderer must integrate $\theta=-\int C\,\mathrm ds$
to match it.

% ===================================================================
\section{Shallow-water equations in $(s,n)$}

\subsection{The starting (Cartesian) system}
Non-rotating depth-averaged shallow water over a bed of still-water depth $H$,
free surface $\eta$, total depth $h=\eta+H$:
\begin{equation}
  \dt\bm u+(\bm u\cdot\nabla)\bm u=-g\nabla\eta+\bm F,
  \qquad
  \dt\eta+\nabla\cdot(h\bm u)=0 .
  \label{eq:sw-cart}
\end{equation}
The pressure is hydrostatic, $p=\rho g(\eta-z)$, so the horizontal pressure
gradient is exactly $-g\nabla\eta$; this is the $L\gg H$ (long-wave) limit.
Keeping $\dt\eta$ is what admits \textbf{gravity waves}; the rigid-lid packages in
\code{../} drop it and are left with the vortical branch only.

\subsection{Curvilinear operators}
For orthogonal coordinates with scale factors $(h_s,h_n)=(\sigma,1)$:
\begin{align}
  \nabla\phi&=\frac{1}{\sigma}\ds\phi\,\hat{\bm t}+\dn\phi\,\hat{\bm n},\\
  \nabla\cdot\bm F&=\frac{1}{\sigma}\Big[\ds F_s+\dn(\sigma F_n)\Big],
  \label{eq:div}\\
  \zeta=(\nabla\times\bm u)\cdot\hat{\bm z}&=\frac{1}{\sigma}\Big[\ds \un-\dn(\sigma \us)\Big].
  \label{eq:vort}
\end{align}
The material acceleration acquires \emph{metric} (centrifugal/Coriolis-like) terms
from the rotation of the basis vectors. With
$\mathrm D/\mathrm Dt=\dt+(\us/\sigma)\ds+\un\dn$ and
$\partial h_s/\partial n=C$, $\partial h_n/\partial s=0$, the standard result is
\begin{equation}
  a_s=\frac{\mathrm D\us}{\mathrm Dt}+\frac{C\,\us\un}{\sigma},
  \qquad
  a_n=\frac{\mathrm D\un}{\mathrm Dt}-\frac{C\,\us^{2}}{\sigma}.
  \label{eq:accel}
\end{equation}

\subsection{The system solved}
Substituting \eqref{eq:div}--\eqref{eq:accel} into \eqref{eq:sw-cart}:
\begin{align}
  &\dt\eta+\frac{1}{\sigma}\Big[\ds(h\us)+\dn(\sigma h\un)\Big]=0,
  \label{eq:cont}\\
  &\dt \us+\frac{\us}{\sigma}\ds \us+\un\dn \us
   +\underbrace{\frac{C\us\un}{\sigma}}_{\text{bend coupling}}
   =-\frac{g}{\sigma}\ds\eta-r_s\us+\nu\nabla^2\us,
  \label{eq:smom}\\
  &\dt \un+\frac{\us}{\sigma}\ds \un+\un\dn \un
   \underbrace{-\frac{C\us^{2}}{\sigma}}_{\text{\bf centrifugal}}
   =-g\,\dn\eta-r_n\un+\nu\nabla^2\un.
  \label{eq:nmom}
\end{align}
The term $-C\us^2/\sigma$ is the whole reason for going to $(s,n)$: it is the
centrifugal acceleration of water rounding a bend, and it is what tilts the free
surface up against the outer bank (\S\ref{sec:base}). No straight-channel model
contains it.

% ===================================================================
\section{Non-dimensionalisation: where the Froude number enters}
\label{sec:nondim}

Scale lengths by the half-width $b$, speeds by the centreline speed
$U_c=U_0+\Delta$, time by $b/U_c$, and depths/$\eta$ by the mean depth $H_0$.
The only place $g$ survives is the pressure gradient, where it appears as
\begin{equation}
  \frac{gH_0}{U_c^{2}}=\frac{1}{F^{2}},\qquad
  \boxed{\;F\equiv\frac{U_c}{\sqrt{gH_0}},\qquad g_{\rm eff}=\frac1{F^{2}}\;}
  \tag{$\to$ \code{g\_eff}}
\end{equation}
so the non-dimensional gravity-wave speed is $\sqrt{g_{\rm eff}H_0}=1/F$.

\paragraph{Why $F$ must be specified, and why one value proves nothing.}
$F$ is an \emph{input} (a property of the river), exactly like a Reynolds number:
without it the system \eqref{eq:cont}--\eqref{eq:nmom} is not defined. At any $F$
\emph{both} families exist --- gravity waves at speed $1/F$ and the vortical
branch at $O(U)$. Which one the bank couples to is decided by \emph{measurement}:
compare $c_{\rm meander}$ with $1/F$, and above all \emph{vary} $F$. If the meander
rode the gravity branch, tripling $1/F$ would move $\sigma,c$. It does not
(\code{sweep\_dispersion.py} $\to$ \code{fig01}), which is the evidence that the
meander is the vortical/Rossby wave.

% ===================================================================
\section{The base state}
\label{sec:base}

\subsection{Parabolic jet and the constant channel-$\beta$}
\begin{equation}
  \boxed{\;\Us(n)=U_0+\Delta\Big(1-\frac{n^{2}}{b^{2}}\Big)\;}
  \tag{$\to$ \code{ubar\_s}}
\end{equation}
so $\Us(0)=U_0+\Delta$ (centre) and $\Us(\pm b)=U_0$ (bank edge), and
\begin{equation}
  \dn\Us=-\frac{2\Delta n}{b^{2}}\ \ (\to\code{ubar\_s\_n}),\qquad
  \dn^{2}\Us=-\frac{2\Delta}{b^{2}}=\text{const}\ \ (\to\code{ubar\_s\_nn}).
\end{equation}
This constant is the load-bearing property. From \eqref{eq:vort} with $\un=0$ and
(for the moment) $C=0$, the base vorticity is $\bar\zeta=-\dn\Us$, hence
\begin{equation}
  \boxed{\;\frac{\partial\bar\zeta}{\partial n}=-\dn^{2}\Us=\frac{2\Delta}{b^{2}}
  =\text{constant}\;}
\end{equation}
--- a uniform cross-channel gradient of background vorticity. This is the exact
analogue of the planetary $\beta$: it is the restoring mechanism that makes a
\emph{vortical (Rossby-type) wave} possible in a river. A parabolic jet is thus not
cosmetic; it is what gives the channel a constant $\beta$. (For $C\neq0$ the same
calculation gives $\bar\zeta=-\dn\Us-C\Us/\sigma$, a curvature correction.)

\subsection{Base curvature = the finite background meander}
\begin{equation}
  \boxed{\;\bar C(s)=\bar C_{\rm amp}\cos(k_m s),\qquad
  \bar C_{\rm amp}=A_{\rm bank}k_m^{2}\ \text{(default)}\;}
  \tag{$\to$ \code{cbar}}
\end{equation}
For a centreline of lateral amplitude $A_{\rm bank}$ and wavenumber $k_m$, the
small-slope curvature is $-\partial_s^2(A_{\rm bank}\cos k_m s)=A_{\rm bank}k_m^{2}\cos k_m s$;
\code{Cbar\_amp} overrides this to set a genuinely tight bend directly.
\textbf{$A_{\rm bank}$ is the fixed background channel shape --- it is \emph{not} the
perturbation seed $A_0$ of \S\ref{sec:seed}, and it does not evolve.}

\subsection{Superelevation}
For a steady base with $\bar U_n=0$ and $\ds(\cdot)=0$, the $n$-momentum
\eqref{eq:nmom} reduces to a balance between the centrifugal term and the pressure
gradient:
\begin{equation}
  \boxed{\;g\,\dn\bar\eta=\frac{\bar C\,\Us^{2}}{\bar\sigma}\;}
  \tag{$\to$ \code{etabar}}
  \label{eq:super}
\end{equation}
integrated as $\bar\eta(s,n)=\dfrac{\bar C(s)}{g}\displaystyle\int_0^n
\frac{\Us^{2}(n')}{1+n'\bar C(s)}\,\mathrm dn'$ (zero on the centreline). The water
surface tilts \emph{up towards the outer bank} --- the superelevation. Note
$\bar C\to0\Rightarrow\bar\eta\to0$ (asserted in the test).

\subsection{Discharge conservation}
The base continuity \eqref{eq:cont} with $\bar U_n=0$ requires
$\ds(\bar h\Us)=0$: the transport is the same through every cross-section. With a
cross-channel-only bed this is automatic ($\Us=\Us(n)$).

\paragraph{What sustains the base.} A parallel jet with bottom drag is not a
steady solution unless something drives it (the downhill slope). As in every
frozen-jet stability analysis we \emph{prescribe} $\Us$ and absorb the residual into
an implicit steady forcing; it cancels from the perturbation equations below.

% ===================================================================
\section{Linearisation}

Write $\us=\Us+\us'$, $\un=\un'$, $\eta=\bar\eta+\eta'$, $h=\bar h+\eta'$,
$C=\bar C+C'$, and drop products of primes.

\subsection{Continuity}
$\ds(h\us)=\ds(\bar h\Us)+\bar h\ds\us'+\Us\ds\eta'$, where the first term vanishes
(discharge), and $\dn(\sigma h\un)\to\dn(\bar\sigma\bar h\,\un')$. Hence
\begin{equation}
  \bar\sigma\,\dt\eta'+\bar h\,\ds\us'+\Us\,\ds\eta'+\dn\!\big(\bar\sigma\bar h\,\un'\big)=0,
  \tag{$\to$ eq.\ 1 in \code{build\_ivp\_SW}}
\end{equation}
already multiplied by $\bar\sigma$ (\S\ref{sec:sigma}).

\subsection{$s$-momentum}
$\frac{\us}{\sigma}\ds\us\to\frac{\Us}{\bar\sigma}\ds\us'$;
$\un\dn\us\to\un'\dn\Us$;
$\frac{C\us\un}{\sigma}\to\frac{\bar C\Us\un'}{\bar\sigma}$. Multiplying by
$\bar\sigma$ collects the two $\un'$ terms into one coefficient:
\begin{equation}
  \bar\sigma\dt\us'+\Us\ds\us'
  +\underbrace{\big(\bar\sigma\,\dn\Us+\bar C\,\Us\big)}_{\to\ \code{coefUn}}\un'
  +g\,\ds\eta'+\bar\sigma r_s\us'-\bar\sigma\nu\nabla^{2}\us'=0 .
  \tag{$\to$ eq.\ 2}
\end{equation}
The $\un'$ coefficient is the cross-channel advection of background momentum ---
proportional to the base vorticity --- and is the term that carries the
Rossby-type restoring.

\subsection{$n$-momentum and the curvature feedback}
\label{sec:feedback}
Expanding $-C\us^{2}/\sigma$ with $C=\bar C+C'$, $\sigma=\bar\sigma+nC'$:
\begin{equation}
  -\frac{C\us^{2}}{\sigma}\;\simeq\;
  \underbrace{-\frac{\bar C\Us^{2}}{\bar\sigma}}_{\text{base}}
  \;\underbrace{-\frac{2\bar C\Us\us'}{\bar\sigma}}_{\to\ \code{twoCbUb}}
  \;\underbrace{-\frac{C'\Us^{2}}{\bar\sigma}}_{\text{\bf new bend}}
  \;+\;\frac{nC'\,\bar C\Us^{2}}{\bar\sigma^{2}} .
\end{equation}
The perturbed pressure term contributes $-\bar\sigma g\dn\eta'-nC'g\dn\bar\eta$, and
using the base balance \eqref{eq:super}, $g\dn\bar\eta=\bar C\Us^2/\bar\sigma$, the
last piece is $-nC'\bar C\Us^{2}/\bar\sigma$. Adding it to the $+nC'\bar C\Us^2/\bar\sigma^2$
above and using $\bar\sigma-n\bar C=1$ leaves the remarkably clean forcing
\begin{equation}
  \boxed{\;\text{$n$-momentum forcing}=\frac{C'\,\Us^{2}}{\bar\sigma}\;\simeq\;C'\Us^{2}.\;}
\end{equation}
So the final equation, with the forcing moved to the left and
$C'=-\partial_s^{2}\zeta_c$ (\S\ref{sec:erosion}):
\begin{equation}
  \bar\sigma\dt\un'+\Us\ds\un'-2\bar C\Us\us'+\bar\sigma g\dn\eta'
  +\bar\sigma r_n\un'-\bar\sigma\nu\nabla^{2}\un'
  +\underbrace{\Us^{2}\,\partial_s^{2}\zeta_c}_{\text{bend $\to$ flow}}=0 .
  \tag{$\to$ eq.\ 3}
\end{equation}
\emph{This is the coupling that closes the meander instability}: a bend creates a
centrifugal forcing, which drives the flow, which erodes the bank, which deepens
the bend.

% ===================================================================
\section{Why every equation is multiplied by $\sigma$}
\label{sec:sigma}

The raw equations carry $1/\sigma=1/(1+nC)$, a \emph{rational} function of $n$.
Dedalus builds such a coefficient as a dense operator (its documentation warns
about exactly this), destroying the banded structure. Multiplying each equation
through by $\sigma$ removes every $1/\sigma$ and leaves coefficients that are
polynomial in $n$ (products of $\sigma$, $\bar h$, $\Us$), so the matrices stay
banded. Multiplying an equation by a nonzero function is an exact,
solution-preserving row operation.
The one residual $1/\bar\sigma$ (in the $C'$ forcing) is dropped at
$O(n\bar C\,C')$, i.e.\ second order in the perturbation amplitude \emph{and} the
base curvature.

% ===================================================================
\section{Bottom friction: the factor of two}

Quadratic (Ch\'ezy) drag is $\bm F=-\dfrac{C_f}{h}|\bm u|\bm u$. With
$\bm u=(\Us+\us',\,\un')$ and $|\bm u|\simeq\Us+\us'$ (the $\un'$ contribution is
second order):
\begin{align}
  |\bm u|\us&\simeq\Us^{2}+2\Us\us'
  &&\Rightarrow\quad r_s=\frac{2C_f\Us}{\bar h}\quad(\to\code{rs\_a}),\\
  |\bm u|\un&\simeq\Us\,\un'
  &&\Rightarrow\quad r_n=\frac{C_f\Us}{\bar h}\quad(\to\code{rn\_a}).
\end{align}
The drag is therefore \textbf{anisotropic}: streamwise perturbations are damped
twice as fast as cross-stream ones, because they perturb $|\bm u|$ as well as the
direction. (The same factor $2$ appears in the ``momentum'' closure of the
rigid-lid packages.) The base part $C_f\Us^2/\bar h$ belongs to the implicit
driving balance and cancels.

% ===================================================================
\section{Boundary conditions, and why there are four tau terms}

The banks sit at fixed $n=\pm b$ in channel coordinates. Water cannot cross a
solid bank, so
\begin{equation}
  \boxed{\;\un'(\pm b)=0\;}\qquad\text{(no penetration)} .
\end{equation}
This assumes \emph{nothing} about the interior divergence: $\nabla\cdot\bm u\neq0$
everywhere inside (that is what the free surface is for). In particular the bank is
\emph{not} treated as a material streamline --- that would be a rigid-lid
(non-divergent) assumption and is inconsistent with \eqref{eq:cont}.

Lateral eddy viscosity $\nu\nabla^{2}$ makes each momentum equation second order in
$n$, which (i) is physically reasonable for a river, (ii) makes the parabolic jet an
actual steady solution, and (iii) gives a standard tau closure. A second-order
operator needs one boundary condition per wall per component, so with free slip
\begin{equation}
  \dn\us'(\pm b)=0\qquad\text{(free slip)} ,
\end{equation}
the tau budget is $2$ per viscous velocity component:
\[
  \underbrace{2}_{\us'}+\underbrace{2}_{\un'}+\underbrace{0}_{\eta'}=4
  \qquad(\to\ \code{t\_us1,t\_us2,t\_un1,t\_un2}),
\]
$\eta'$ needing none because its equation contains no second $n$-derivative.
Counting unknowns and equations:
$\{\us',\un',\eta',\zeta_c\}+4\ \text{tau}=8$ and
$3\ \text{PDE}+2\ \text{no-penetration}+2\ \text{free-slip}+1\ \text{erosion}=8$.
The tau terms are lifted onto the last two Chebyshev modes,
\code{lift(t,-1)+lift(t,-2)} on \code{ybasis.derivative\_basis(2)}.

% ===================================================================
\section{Erodible bank: the meander equation}
\label{sec:erosion}

The bank is a sediment boundary, not a material line: it recedes because the flow
erodes it. Ikeda's closure makes the migration rate proportional to the
\emph{near-bank excess longitudinal velocity}. For a \textbf{bend} the outer bank is
fast and the inner bank slow, so it is the \textbf{antisymmetric} part that moves the
centreline:
\begin{equation}
  \boxed{\;\dt\zeta_c=E\cdot\tfrac12\big[\us'(+b)-\us'(-b)\big],\qquad E=\varepsilon\,U_0\;}
  \tag{$\to$ eq.\ 8, \code{bank\_E}}
\end{equation}
where $\zeta_c(s,t)$ is the lateral offset of the centreline (constant width).

\paragraph{Why antisymmetric.} Using two independent banks with
$\dt\xi_\pm=E\us'(\pm b)$ instead makes the centreline
$\tfrac12(\xi_++\xi_-)$ respond to the \emph{symmetric} part of the near-bank
velocity. That is the channel-\emph{widening} mode, not bending: in practice it
gives $\sigma<0$ (no meander instability). Parity matters.

\paragraph{Closing the loop.} The migrating centreline changes the curvature,
\begin{equation}
  C'(s,t)=-\partial_s^{2}\zeta_c ,
\end{equation}
which re-enters the $n$-momentum as the forcing derived in \S\ref{sec:feedback}.
The cycle
\[
  \zeta_c\ \longrightarrow\ C'\ \longrightarrow\ \text{flow}\ \longrightarrow\
  \us'(\pm b)\ \longrightarrow\ \dt\zeta_c
\]
is the meander instability. Its growth rate is what the IVP measures.

% ===================================================================
\section{Seeding, and why the seed must be a domain mode}
\label{sec:seed}

The base state is an exact steady solution: started there, nothing ever happens.
Stability analysis therefore \emph{requires} a perturbation, and its growth rate is
the answer. We seed a pure bend with zero initial flow:
\begin{equation}
  \zeta_c(s,0)=A_0\cos(k_{\rm eff}s),\qquad \us'=\un'=\eta'=0 .
\end{equation}
\paragraph{Mode snapping.} $\cos(ks)$ is periodic on $[0,L_s]$ only if $kL_s/2\pi$
is an integer. Otherwise the seed has a jump at the domain edge, whose spectral
content spreads over many modes and grows as a \emph{localized packet} instead of a
clean single-wavelength meander. We therefore snap to the nearest exact mode:
\begin{equation}
  m=\mathrm{round}\!\Big(\frac{kL_s}{2\pi}\Big),\qquad
  k_{\rm eff}=\frac{2\pi m}{L_s} .
\end{equation}
\paragraph{$A_0$ is a free constant.} The system is linear, so if
$(\us',\un',\eta',\zeta_c)$ solves it then so does $\lambda\times$ it. Only
$\sigma$, $c$ and the \emph{ratios} between fields are physical; the overall size is
set by whatever initial disturbance one assumes. This is why the movie may fix the
overall constant once (final meander $=0.5b$) and apply it to every frame and every
field --- that is a re-choice of $A_0$, not a distortion.

% ===================================================================
\section{Measuring $\sigma$ and $c$}

The centreline is demodulated onto its seeded mode,
$a(t)=\widehat{\zeta_c}(m,t)$, and
\begin{equation}
  \sigma=\frac{\mathrm d}{\mathrm dt}\ln|a|,\qquad
  c=-\frac1k\frac{\mathrm d}{\mathrm dt}\arg a ,
  \tag{$\to$ \code{measure\_sigma\_c}}
\end{equation}
both fitted over the last two thirds of the run (dropping the initial adjustment,
during which the flow spins up from rest). $\sigma>0$ is an unstable meander;
$c>0$ means the bend train migrates downstream.

% ===================================================================
\section{What the model deliberately does \emph{not} contain}

Strictly linear perturbation dynamics (no $\bm u'\!\cdot\!\nabla\bm u'$); no
sediment (Exner) feedback --- the bed $H$ is prescribed, only the banks move; no
rotation ($f=0$, a river); no vertical structure, hence no helical secondary flow
(it is depth-averaged away); and the base state is frozen (its steady residual is an
implicit forcing). The lateral viscosity uses the flat Laplacian rather than the
exact curvilinear vector Laplacian --- it is a wall-closure and short-wave
regulariser, small, and not part of the physics being measured.

\end{document}
